% GENERATED FILE. Edit canonical lessons, then run npm run generate:books.
% canonical-source-hash: fnv1a64:ed9831e7582edd37
% canonical-lessons: RU-C483-tsarapina, RU-C483-ushibitsya, RU-C483-porez, RU-C483-ssadina, RU-C483-zazhivat

\chapter{Scratches and Grazes}
\label{ch:ru-scratches-and-grazes}
\hlchaptermodality{483}

\begin{chapteropening}
\textbf{By the end of this chapter:} \emph{I can talk about small injuries and how they heal.}

Complete the last lesson of chapter 483: I can talk about small injuries and how they heal.
\end{chapteropening}

\section[\texorpdfstring{\ru{царапина} (tsarápina)}{царапина (tsarápina)}]{\ru{царапина} (tsarápina) — a scratch}

\label{lesson:RU-C483-tsarapina}



\begin{quote}
Before the new one: say the Russian for a nostril, then the Russian for a jaw.
\end{quote}

\subsection*{You'll want to know: \ru{царапина}}
\textbf{\ru{царапина}} (\emph{tsarápina}) — \textquotedblleft{}a scratch\textquotedblright{}.

\begin{quote}
\ru{После} \ru{игры} \ru{с} \ru{кошкой} \ru{у} \ru{меня} \ru{на} \ru{руке} \ru{царапина}.
\end{quote}

\begin{quote}
\emph{Pósle igrý s kóshkoy u menyá na ruké tsarápina.}
\end{quote}

\begin{quote}
After playing with the cat I've got a scratch on my hand.
\end{quote}

Feminine. The verb is \emph{\ru{царапать} / \ru{поцарапать}} — to scratch.

\begin{cousinweb}
\emph{\ru{царапина}} is built on \emph{\ru{царапать}} 'to scratch', which EtymDB lists without a source; it is usually taken as imitative.
\end{cousinweb}

\begin{grammarlens}[title={What you've built}]
One more word for telling, explaining and describing at B1.
\end{grammarlens}

\subsection*{Your turn}
\begin{itemize}
\raggedright
  \item \emph{Say it:} \emph{tsarápina}
  \item \emph{Say it:} \emph{tsarápina}, once more
  \item \emph{Say it:} a sentence of your own with \emph{tsarápina}
  \item \emph{Recall:} say the Russian for a nostril, then the Russian for a jaw, then say \emph{tsarápina} again
\end{itemize}

\begin{tcolorbox}[breakable,colback=teal!4,colframe=teal!35!black,title={Before you move on}]
What does \ru{царапина} mean? (\textquotedblleft{}A scratch\textquotedblright{}.) Say it once more.
\end{tcolorbox}
\section[\texorpdfstring{\ru{ушибиться} (ushibítsya)}{ушибиться (ushibítsya)}]{\ru{ушибиться} (ushibítsya) — to hurt oneself (by a knock)}

\label{lesson:RU-C483-ushibitsya}



\begin{quote}
Before the new one: say the Russian for a jaw, then the Russian for a scratch.
\end{quote}

\subsection*{You'll want to know: \ru{ушибиться}}
\textbf{\ru{ушибиться}} (\emph{ushibítsya}) — \textquotedblleft{}to hurt oneself (by a knock)\textquotedblright{}.

\begin{quote}
\ru{Ребёнок} \ru{упал} \ru{и} \ru{сильно} \ru{ушибился}.
\end{quote}

\begin{quote}
\emph{Rebyónok upál i sílno ushibílsya.}
\end{quote}

\begin{quote}
The child fell and hurt himself badly.
\end{quote}

Perfective; past \ru{ушибся}, \ru{ушиблась}. Imperfective \ru{ушибаться}. The noun is \emph{\ru{ушиб}} — a bruise.

\begin{cousinweb}
\emph{\ru{ушибиться}} is the reflexive of \emph{\ru{ушибить}} 'to bruise', built on \emph{\ru{у}-} and \emph{\ru{шибить}} 'to strike'; EtymDB lists \emph{\ru{ушибить}} without a source and has no entry for \emph{\ru{шибить}}.
\end{cousinweb}

\begin{grammarlens}[title={What you've built}]
One more word for telling, explaining and describing at B1.
\end{grammarlens}

\subsection*{Your turn}
\begin{itemize}
\raggedright
  \item \emph{Say it:} \emph{ushibítsya}
  \item \emph{Say it:} \emph{ushibítsya}, once more
  \item \emph{Say it:} a sentence of your own with \emph{ushibítsya}
  \item \emph{Recall:} say the Russian for a jaw, then the Russian for a scratch, then say \emph{ushibítsya} again
\end{itemize}

\begin{tcolorbox}[breakable,colback=teal!4,colframe=teal!35!black,title={Before you move on}]
What does \ru{ушибиться} mean? (\textquotedblleft{}To hurt oneself (by a knock)\textquotedblright{}.) Say it once more.
\end{tcolorbox}
\section[\texorpdfstring{\ru{порез} (poréz)}{порез (poréz)}]{\ru{порез} (poréz) — a cut}

\label{lesson:RU-C483-porez}



\begin{quote}
Before the new one: say the Russian for a scratch, then the Russian for to hurt oneself.
\end{quote}

\subsection*{You'll want to know: \ru{порез}}
\textbf{\ru{порез}} (\emph{poréz}) — \textquotedblleft{}a cut\textquotedblright{}.

\begin{quote}
\ru{Мама} \ru{промыла} \ru{порез} \ru{и} \ru{заклеила} \ru{его} \ru{пластырем}.
\end{quote}

\begin{quote}
\emph{Máma promýla poréz i zakléila yegó plástyrem.}
\end{quote}

\begin{quote}
Mum washed the cut and put a plaster on it.
\end{quote}

Masculine.

\begin{cousinweb}
\emph{\ru{порез}} is built on \emph{\ru{по}-} and the root of \emph{\ru{резать}} 'to cut', which is inherited from Old East Slavic and Proto-Slavic \emph{rězati} (EtymDB).
\end{cousinweb}

\begin{grammarlens}[title={What you've built}]
One more word for telling, explaining and describing at B1.
\end{grammarlens}

\subsection*{Your turn}
\begin{itemize}
\raggedright
  \item \emph{Say it:} \emph{poréz}
  \item \emph{Say it:} \emph{poréz}, once more
  \item \emph{Say it:} a sentence of your own with \emph{poréz}
  \item \emph{Recall:} say the Russian for a scratch, then the Russian for to hurt oneself, then say \emph{poréz} again
\end{itemize}

\begin{tcolorbox}[breakable,colback=teal!4,colframe=teal!35!black,title={Before you move on}]
What does \ru{порез} mean? (\textquotedblleft{}A cut\textquotedblright{}.) Say it once more.
\end{tcolorbox}
\section[\texorpdfstring{\ru{ссадина} (ssádina)}{ссадина (ssádina)}]{\ru{ссадина} (ssádina) — a graze}

\label{lesson:RU-C483-ssadina}



\begin{quote}
Before the new one: say the Russian for to hurt oneself, then the Russian for a cut.
\end{quote}

\subsection*{You'll want to know: \ru{ссадина}}
\textbf{\ru{ссадина}} (\emph{ssádina}) — \textquotedblleft{}a graze\textquotedblright{}.

\begin{quote}
\ru{У} \ru{него} \ru{на} \ru{колене} \ru{большая} \ru{ссадина}.
\end{quote}

\begin{quote}
\emph{U negó na koléne bolsháya ssádina.}
\end{quote}

\begin{quote}
He has a big graze on his knee.
\end{quote}

Feminine.

\begin{cousinweb}
\emph{\ru{ссадина}} is usually linked to \emph{\ru{ссадить}} in the sense 'to scrape off skin'; EtymDB has no entry, so the link is not settled.
\end{cousinweb}

\begin{grammarlens}[title={What you've built}]
One more word for telling, explaining and describing at B1.
\end{grammarlens}

\subsection*{Your turn}
\begin{itemize}
\raggedright
  \item \emph{Say it:} \emph{ssádina}
  \item \emph{Say it:} \emph{ssádina}, once more
  \item \emph{Say it:} a sentence of your own with \emph{ssádina}
  \item \emph{Recall:} say the Russian for to hurt oneself, then the Russian for a cut, then say \emph{ssádina} again
\end{itemize}

\begin{tcolorbox}[breakable,colback=teal!4,colframe=teal!35!black,title={Before you move on}]
What does \ru{ссадина} mean? (\textquotedblleft{}A graze\textquotedblright{}.) Say it once more.
\end{tcolorbox}
\section[\texorpdfstring{\ru{заживать} (zazhivát)}{заживать (zazhivát)}]{\ru{заживать} (zazhivát) — to heal (of a wound)}

\label{lesson:RU-C483-zazhivat}



\begin{quote}
Before the new one: say the Russian for a cut, then the Russian for a graze.
\end{quote}

\subsection*{You'll want to know: \ru{заживать}}
\textbf{\ru{заживать}} (\emph{zazhivát}) — \textquotedblleft{}to heal (of a wound)\textquotedblright{}.

\begin{quote}
\ru{Рана} \ru{заживает} \ru{очень} \ru{медленно}.
\end{quote}

\begin{quote}
\emph{Rána zazhiváyet óchen médlenno.}
\end{quote}

\begin{quote}
The wound is healing very slowly.
\end{quote}

Imperfective; perfective \ru{зажить} (\ru{заживёт}). Of a wound or a cut: \emph{\ru{рана} \ru{заживает}}.

\begin{cousinweb}
\emph{\ru{заживать}} is built on \emph{\ru{за}-} and \emph{\ru{жить}} 'to live', which is inherited from Proto-Slavic \emph{žiti} (EtymDB): a wound that heals comes back to life.
\end{cousinweb}

\begin{grammarlens}[title={What you've built}]
One more word for telling, explaining and describing at B1.
\end{grammarlens}

\subsection*{Your turn}
\begin{itemize}
\raggedright
  \item \emph{Say it:} \emph{zazhivát}
  \item \emph{Say it:} \emph{zazhivát}, once more
  \item \emph{Say it:} a sentence of your own with \emph{zazhivát}
  \item \emph{Recall:} say the Russian for a cut, then the Russian for a graze, then say \emph{zazhivát} again
\end{itemize}

\begin{tcolorbox}[breakable,colback=teal!4,colframe=teal!35!black,title={Before you move on}]
What does \ru{заживать} mean? (\textquotedblleft{}To heal (of a wound)\textquotedblright{}.) Say it once more.
\end{tcolorbox}
